\section{Method}\label{sec:method}
The APEIRON pipeline consists of three foundational layers that operate identically regardless of the input modality.

\subsection{Layer~1: Ultimate Observables}
For a real‑valued time series $x(t)$, we partition the signal into overlapping windows of length $W$ with stride $S$.
Each window is a vector $\mathbf{w}_i \in \mathbb{R}^W$ and is instantiated as an Ultimate Observable with observability type \texttt{CONTINUOUS}.
\footnote{While a window is technically decomposable at the digital‑sampling level, it is treated as an \emph{operational atom} in the category $\mathcal{C}_1$ relative to the emergence of Layer~2 structures. Atomicity in APEIRON~\cite{Beniers2026Categorical} is scale‑dependent: the window is the minimal unit that yields a meaningful co‑activation signal for the construction of the relational hypergraph.}
No normalisation beyond $\ell_2$ unit‑norm is applied.

\subsection{Layer~2: Relational Hypergraph}
The co‑activation strength between two observables $\mathbf{w}_i$ and $\mathbf{w}_j$ is defined as the absolute Pearson correlation of their window vectors:
\[
A_{ij} = \left|\mathrm{corr}(\mathbf{w}_i, \mathbf{w}_j)\right|.
\]
The diagonal is set to zero.
The resulting weighted adjacency matrix defines the relational hypergraph of the system.

\subsection{Layer~3: Functional Clustering}
Louvain community detection~\cite{Blondel2008} is applied to the adjacency matrix to partition the windows into functional clusters $\mathcal{C}_1, \dots, \mathcal{C}_k$.
No pre‑specification of the number of clusters is required.

\subsection{Functional validation: counterfactual ablation}
To verify that a discovered cluster carries genuine functional information, we define a \emph{Functional Significance Score} $S_k$ for cluster $\mathcal{C}_k$.
Let $f$ be a logistic regression classifier trained on the one‑hot encoded cluster membership vectors, and let $\mathrm{Acc}(f,X)$ denote the classification accuracy of $f$ on a test set $X$.
The score is defined as
\[
S_k = \mathrm{Acc}\bigl(f, X\bigr) - \mathrm{Acc}\bigl(f, X \setminus \mathcal{C}_k\bigr),
\]
where $X \setminus \mathcal{C}_k$ denotes the test set with the columns corresponding to cluster $\mathcal{C}_k$ set to zero (\emph{feature ablation}), while the number of test samples remains unchanged.
A large, statistically significant value of $S_k$ indicates that the cluster encodes functionally relevant structure that the classifier exploits.